\begin{center}
\textsc{Gauss} an \textsc{Gerling}.\quad Göttingen, 21.~Februar 1822.
\end{center}

\par\noindent\hspace*{\parindent}\mbox{.\ .\ .\ .\ .\ .}\quad Wenn der Anschluss an die Göttinger Sternwarte und der an den Inselsberg nicht von einem einzigen Platz auf dem Meisner aus erfolgen können, dürfen Sie nun \textls[120]{beide} Plätze nutzen.
Der zusätzliche Aufwand ist dabei sehr gering, denn nach den Prinzipien der Wahrscheinlichkeitsrechnung brauchen Sie, sofern Sie an beiden Punkten messen, die Winkel zwischen den von beiden Punkten aus sichtbaren Objekten an jedem Punkt nur halb so oft zu wiederholen, wie Sie es bei nur einem Standpunkt getan hätten.
Die gegenseitige Lage der beiden Punkte mit größter Genauigkeit und doch ohne unnötige Mühe zu bestimmen, wird Ihnen nicht schwer \mbox{fallen.\quad.\ .\ .\ .\ .\ .}
